\documentclass{article}
\usepackage{iclr2026_conference,times}
\input{math_commands.tex}
\usepackage{graphicx}
\usepackage{booktabs}
\usepackage{hyperref}
\usepackage{url}
\usepackage{float}

\title{Encoder-Preserving Mixed Precision Quantization Improves DINO-WM Planning on Wall}

\author{}

\begin{document}
\maketitle

\begin{abstract}
We study whether preserving encoder precision improves planning quality under strict efficiency constraints in world models.
Using DINO-WM on the Wall environment, we compare FP16, uniform INT8 weight-only quantization, and mixed INT8 quantization that keeps the encoder in FP16.
In a Mac-only preliminary setting, mixed precision improves planning success over uniform INT8 while retaining a similar local efficiency footprint.
We release a replay-based demo and a reproducible script bundle for all runs.
\end{abstract}

\section{Introduction}
Efficient spatial reasoning requires planning systems that remain accurate under constrained memory and latency budgets.
World models are attractive because planning can be performed in latent space, but precision reduction can degrade representation quality.
This work focuses on a practical mixed precision strategy that preserves encoder precision while quantizing predictive dynamics.

\section{Setup}
We evaluate DINO-WM \citep{zhou2024dinowm} on the Wall task using the repository's native planning pipeline and success metric.
Success is parsed directly from planning output logs, without redefining task metrics.
All variants use identical checkpoints, seeds, and planning settings.

\section{Method}
We use weight-only INT8 quantization for linear layers.
We compare three variants: (1) FP16 baseline, (2) uniform INT8 on encoder and predictor, and (3) mixed INT8 where the encoder remains FP16 and only predictor layers are quantized.
The hypothesis is that preserving encoder precision better maintains goal-relevant latent structure.

\section{Experiments}
We run Wall-only evaluation with planning step budgets and multiple seeds.
Each run records success rate, average planning time, peak memory proxy, and model size.

\begin{figure}[H]
  \centering
  \IfFileExists{../figures/success_vs_opt_steps.pdf}{
    \includegraphics[width=0.95\linewidth]{../figures/success_vs_opt_steps.pdf}
  }{
    \fbox{\parbox[c][4.0cm][c]{0.95\linewidth}{\centering Placeholder for success-vs-opt-steps figure}}
  }
  \caption{Planning success versus optimization steps for FP16, uniform INT8, and mixed INT8.}
  \label{fig:success}
\end{figure}

\begin{figure}[H]
  \centering
  \IfFileExists{../figures/efficiency_table.pdf}{
    \includegraphics[width=0.95\linewidth]{../figures/efficiency_table.pdf}
  }{
    \fbox{\parbox[c][2.5cm][c]{0.95\linewidth}{\centering Placeholder for efficiency table}}
  }
  \caption{Efficiency comparison across variants (model size, planning latency, and peak GPU memory).}
  \label{fig:eff}
\end{figure}

\section{Demo}
We provide a replay-based Streamlit viewer that loads precomputed episode trajectories and allows instant variant toggling on a Mac without GPU dependencies.

\section{Limitations}
This MVP study is limited to one environment (Wall), one quantization family (weight-only INT8), and no activation quantization.
For Mac-only execution we use a fake INT8 weight-only fallback instead of CUDA kernel-backed int8 inference, so efficiency numbers are preliminary local proxies.
Broader environment coverage and richer quantization schemes remain future work.

\bibliography{references}
\bibliographystyle{iclr2026_conference}

\end{document}
